\documentclass[11pt]{article}
\usepackage[letterpaper,margin=1.02in]{geometry}
\usepackage[T1]{fontenc}
\usepackage{lmodern,microtype,amsmath,amssymb,amsthm,mathtools,booktabs,xcolor,hyperref}
\definecolor{ink}{HTML}{155E69}
\hypersetup{colorlinks=true,linkcolor=ink,urlcolor=ink,citecolor=ink,pdftitle={The exact cost of matching Dirichlet zero ordinates}}
\newtheorem{theorem}{Theorem}
\newtheorem{lemma}[theorem]{Lemma}
\newtheorem{corollary}[theorem]{Corollary}
\newcommand{\E}{\mathbb E}
\newcommand{\R}{\mathbb R}
\newcommand{\var}{\operatorname{Var}}
\newcommand{\TV}{\operatorname{TV}}
\newcommand{\cost}{\mathcal C}
\newcommand{\A}{\mathcal A}
\newcommand{\B}{\mathcal B}
\newcommand{\norm}[1]{\left\lVert #1\right\rVert}
\newcommand{\abs}[1]{\left\lvert #1\right\rvert}
\title{The exact cost of matching Dirichlet zero ordinates}
\author{Research note for the ProveIt project}
\date{27 November 2026}
\begin{document}
\maketitle
\begin{abstract}
For two distinct primitive Dirichlet characters, list the ordinates of the nontrivial zeros at heights between $T$ and $2T$. Pairing two entries costs their distance. Leaving an entry unmatched costs $\lambda$.
If $0<\lambda\le1$ and $M=\lambda\log T\longrightarrow\infty$, the minimum total cost is asymptotic to
$\sqrt{2}\,T\sqrt{\log M}/\pi^{3/2}$. For equal-degree products the constant is the Euclidean distance of their multiplicity vectors divided by $\pi^{3/2}$. When every occurrence of the smaller list must be matched, we identify the additional dependence on the two endpoint counts. For equal product conductors the random-height limiting law has an atom of mass one half. The proof uses prime orthogonality and the dispersion of the zero-count difference about its local median. All assertions for ordinates are unconditional.
\end{abstract}

\section{The extremal problem}
Let $\chi_1,\ldots,\chi_d$ be a fixed family of distinct primitive Dirichlet characters of conductors $q_j$, allowing the character of conductor one. For $a,b\in\mathbb Z_{\ge0}^d$ write $c=a-b$. Throughout,
\[
 \sum_j a_j=\sum_j b_j,\qquad c\ne0,\qquad
 \kappa_c=\frac{\norm{c}_2}{\pi^{3/2}},\qquad
 \beta_c=\frac{1}{2\pi}\sum_j c_j\log q_j.
\]
Let $\A,\B$ be the multisets of imaginary parts of the nontrivial zeros in $(T,2T]$ of $\prod_jL(s,\chi_j)^{a_j}$ and $\prod_jL(s,\chi_j)^{b_j}$, respectively. Repeated zeros are counted repeatedly. A partial matching $P$ uses each occurrence at most once. Define
\[
 D(P)=|\A|+|\B|-2|P|,\qquad
 \cost_\lambda=\min_P\left\{\sum_{(x,y)\in P}|x-y|+\lambda D(P)\right\}.
\]

\begin{theorem}[Matching with a deletion penalty]\label{th:main}
For every $0<\lambda=\lambda(T)\le1$ with $M=\lambda\log T\to\infty$,
\[
 \cost_\lambda=(\kappa_c+o(1))\,T\sqrt{\log M}.
\]
The quantifier includes arbitrarily slowly diverging $M$. The family and the multiplicity vectors are fixed.
\end{theorem}

\begin{theorem}[Matching the entire smaller list]\label{th:full}
If $\beta_c\ne0$, let $W_T$ be the minimum total distance with exactly $\min(|\A|,|\B|)$ pairs. Then
\[
 W_T=(\kappa_c+o(1))\,T\sqrt{\log\log T}.
\]
In fact, uniformly for $\lambda\ge1$,
\[
 \cost_\lambda=\lambda\bigl||\A|-|\B|\bigr|
       +(\kappa_c+o(1))\,T\sqrt{\log\log T}.
\]
\end{theorem}

\begin{theorem}[Equal product conductors]\label{th:end}
Suppose $\beta_c=0$, put $s_T=\sqrt{\log\log T}$ and write $F_0(t)=\sum_j c_j N_j(t)$, where $N_j$ counts nontrivial positive ordinates through $t$. Let $C_c$ be the constant smooth term in (3). Let $z_T$ be the nearest point to zero in the closed interval with endpoints
\[
  (F_0(T)-C_c)/s_T,\qquad (F_0(2T)-C_c)/s_T .
\]
For $G\sim\mathcal N(0,\norm{c}_2^2/(2\pi^2))$ the maximum-cardinality matching obeys
\[
  \frac{W_T}{T s_T}=\E|G-z_T|+o(1).
\]
The error is additive, including along sequences with large endpoint discrepancies. A bounded change of the centering constant changes the right side by $O(1/s_T)$.
\end{theorem}

For example, for $\zeta(s)$ and the Dirichlet beta function $L(s,\chi_{-4})$, the second list has $(\log4)/(2\pi)\,T+O(\log T)$ extra occurrences. Even with the best choice of these omissions, the mean displacement per matched zero is asymptotic to
$2\sqrt{2/\pi}\,\sqrt{\log\log T}/\log T$. Thus a mismatch of the smooth densities and fluctuations about those densities enter in different orders.

\paragraph{Connection with the project.}
The merged ProveIt article \cite{repo}, section on recovering zeros and characters, proves lower bounds for moving or deleting zeros in short windows. It also introduces the cost of one per omission and $\min(1,|x-y|)$ per pair. On our dyadic lists this metric is exactly $\cost_{1/2}+\tfrac12\bigl||\A|-|\B|\bigr|$. Indeed one may pair all occurrences of the smaller list after capping the movement at one, and in $\cost_{1/2}$ an overlong pair may instead be omitted at cost one. Theorem~\ref{th:main} gives the leading term of the repository's metric in this height regime, since the cardinality difference is $O(T)$. For any finite family of equal-degree products, the entire matrix of these distances, divided by $T\sqrt{\log\log T}$, converges to the Euclidean distance matrix of the exponent vectors divided by $\pi^{3/2}$.

Fujii \cite{fujii} studied separation and ordering for zeros of distinct Dirichlet functions. Bombieri and Perelli \cite{bp} obtained positive-proportion noncoincidence using a joint short-increment limit theorem. Bourgade \cite{bourgade} studied correlations of mesoscopic zero counts at shared endpoints. We use the prime approximation behind these results to solve a distance minimization problem. The proof counts ordinates in the whole strip. Under the project's Dirichlet GRH input, the same distances are the Euclidean distances of the zeros in the complex plane. Neither simplicity nor joint pair correlation enters.

\section{Two inequalities for finite lists}
For real multisets on $(l,r]$, set $F(x)=\#(\A\cap(l,x])-\#(\B\cap(l,x])$, with multiplicities understood. Adding a constant to $F$ below is allowed provided the boundary values are translated as well. For an interval $I$ define
\[
  H_I(F)=\inf_{m\in\R}\int_I |F(x)-m|\,dx .
\]
Any median for uniform $x\in I$ attains the infimum. For real functions $F,G$ we have
$|H_I(F)-H_I(G)|\le\int_I |F-G|$. Also $H_I$ is unchanged by adding a constant.

\begin{lemma}[Local dispersion and a transport upper bound]\label{le:flow}
Let $J=F(r)-F(l)$. For any disjoint intervals $I$ of length $h\le\lambda$ inside $[l,r]$,
\[
 \cost_\lambda\ \ge\ (\lambda-h)|J|+\sum_I H_I(F).                 \tag{1}
\]
For any integer-valued step function $g$ on $(l,r)$ with finitely many jumps,
\[
 \cost_\lambda\ \le\ \int_l^r|F-g|\,dx+\lambda V(g),             \tag{2}
\]
where $V(g)$ is its total variation including the jumps from $F(l)$ on the left and to $F(r)$ on the right. If the extended $g$ is monotone then $W_T\le\int_l^r|F-g|\,dx$ for the analogous maximum-cardinality matching of these lists.
\end{lemma}

\begin{proof}
Any $1$-Lipschitz $f$ with $|f|\le\lambda$ gives the lower bound $\sum_{\A}f-\sum_{\B}f$ by summing over a matching. In one block choose a median $m$. Take $f_I'$ equal to $-\operatorname{sgn}(F-m)$ away from the median level and assign values in $[-1,1]$ on that level to give mean zero. This is possible by the median inequalities. Its primitive $f_I$, zero at both ends, extended by zero, satisfies $\int f_I\,dF=H_I(F)$ by integration by parts and $|f_I|\le h$. Adding the block primitives to $(\lambda-h)\operatorname{sgn}J$ gives an admissible test and proves (1).

For (2) form a path on all list entries, jump positions and endpoints, and add one auxiliary vertex linked to each possible jump at cost $\lambda$ per unit. Send signed integral flow $F-g$ on each consecutive interval, at cost its length times the absolute flow. Send the signed jumps of $g$ to the auxiliary vertex, including the end corrections. Flow conservation gives the prescribed excess of first-list over second-list entries at each location and the net balancing amount at the auxiliary vertex. Decompose integral flow into unit source-to-sink paths after discarding cycles. A path between two list entries which avoids the auxiliary vertex can match the two for at most its cost. A path through that vertex can pay for the two omissions, and a path terminating or starting there pays for one. Cancel coincident entries at zero cost first. This proves (2). If $g$ is monotone all auxiliary edges carrying flow have the same orientation. Hence exactly $|J|$ occurrences are omitted and the remaining transport costs at most the horizontal part of the flow.
\end{proof}

\section{The prime frequencies on one block}
Write $N_j(x)$ for the nontrivial zero count up to positive height $x$ and
$F_0(x)=\sum c_jN_j(x)$. Local dispersions agree for $F_0$ and $F=F_0-F_0(T)$.
The one-sided Riemann--von Mangoldt formula for a primitive character gives, away from the countable exceptional heights,
\[
 F_0(x)=m(x)+\sum_j\frac{c_j}{\pi}\Im\log L(1/2+ix,\chi_j)+O(1/x),
 \qquad m(x)=\beta_c x+C_c .                                         \tag{3}
\]
The logarithm is continued horizontally from the right. The common term $(x/2\pi)\log(x/(2\pi e))$ cancels. In particular $F_0(x)=m(x)+O(\log(2+x))$ also at endpoints under our counting convention. The standard argument estimate allows the full jump at a zero.

Consider $0<h\le1$, $M=h\log T\to\infty$, and set
\[
 L=\log M,\quad X=\exp\{\log T/L^{1/4}\},\quad Y=\exp(1/h).
\]
Eventually $Y<X$, $X=T^{o(1)}$ and $\log(\log X/\log Y)=L-\tfrac14\log L$. Put $a_p=\sum_j c_j\chi_j(p)$.
Let $w(s)=v(e^s)$ be the smooth weight of \cite[\S2]{bp}, bounded by one, with $w(0)=1$ and $w(s)=0$ for $s\ge1$. Split the prime sum into
\[
 P_{\rm lo}(x)=\frac1\pi\Im\sum_{p\le Y}\frac{a_p w(\log p/\log X)}{\sqrt p}\,p^{-ix},
 \qquad
 P_{\rm hi}(x)=\frac1\pi\Im\sum_{Y<p\le X}\frac{a_p w(\log p/\log X)}{\sqrt p}\,p^{-ix}.
\]
In this section averages in $x$ are uniform over $[T,2T]$, written $\E_x$. All estimates used with a shift of absolute size at most two hold over the shifted interval as well.

\begin{lemma}[Prime approximation and local Gaussian values]\label{le:prime}
With $B=m+P_{\rm lo}$ and $Z=P_{\rm hi}$, the following hold.
\setcounter{equation}{3}
\begin{align}
 \E_x|F_0(x)-B(x)-Z(x)|&=O(1+L^{1/4}),\label{eq:err}\\
 \sup_{|v|\le h}\E_x|B(x+v)-B(x)|&=O(1).\label{eq:slow}
\end{align}
For fixed distinct $u_1,\ldots,u_s\in[0,1]$,
\[
 (Z(x+hu_1),\ldots,Z(x+hu_s))/\sqrt L
 \ \Longrightarrow\ (G_1,\ldots,G_s),\qquad
 G_j\text{ independent }\mathcal N(0,\sigma_c^2),\quad
 \sigma_c^2=\frac{\norm{c}_2^2}{2\pi^2}.
\]
The normalized second moments of the individual coordinates are bounded uniformly in their offsets.
\end{lemma}

\begin{proof}
For the fixed primitive characters, hypotheses B and D of \cite{bp} hold unconditionally. B follows from the Euler product, functional equation and the prime number theorem in fixed arithmetic progressions, and D is the Dirichlet zero-density estimate recorded on p.~273 of that paper. Their Lemma~1 approximates each logarithm in mean absolute value by its smoothed prime-power sum with error $O(\log T/\log X)$. The terms $p^k$ with $k\ge2$ cost $O(1)$ in mean square, since their squared coefficients sum to $O(\sum_p\sum_{k\ge2}p^{-k})$ and the Dirichlet polynomial has length $X=o(T)$. The mean-value theorem for Dirichlet polynomials gives the assertion, whence \eqref{eq:err}. For shifted intervals the same nonnegative integral estimates follow by covering with a bounded number of dyadic intervals of comparable height. For \eqref{eq:slow} use the same mean-value theorem and
\[
 \sum_{p\le Y}\frac1p |p^{-iv}-1|^2
 \ll h^2\sum_{p\le Y}\frac{(\log p)^2}{p}\ll h^2(\log Y)^2=1.
\]
The affine main term moves by $O(h)$.

For clarity we give the frequency calculation for the joint assertion. By orthogonality of fixed distinct primitive characters, for $y\ge3$ and any fixed $A>0$,
\[
 \sum_{p\le y}\frac{|a_p|^2}{p}
 =\norm{c}_2^2\log\log y+C+O_A((\log y)^{-A}).                     \tag{6}
\]
Indeed two distinct primitive characters induce distinct characters on the units modulo the least common multiple of their conductors, so each cross term has mean zero on primes in reduced residue classes. The prime number theorem in these fixed progressions with its classical error and partial summation yield (6). In sums of $1/p$ up to $X$ removing the squared weight costs $O(1)$ in absolute value, since $|1-w(s)^2|\ll s$ for $0\le s\le1$. Hence the diagonal high-frequency sum is $\norm{c}_2^2L+o(L)$. For fixed $0<|u|\le1$,
\[
 \sum_{Y<p\le X}\frac{|a_p|^2}{p} w(\log p/\log X)^2
      \cos(uh\log p)=O_u(1).                                         \tag{7}
\]
The main integral is bounded by changing variables to $v=h\log p$ and integrating $\cos(uv)/v$ from $1$ to $h\log X$. The error follows by integration by parts in (6) with $A=2$, as $h\le1$ and $\log Y\ge1$.

For real $b_1,\ldots,b_s$, the corresponding linear combination is the imaginary part of a prime polynomial with coefficients
\[
 d_p=\frac{a_pw(\log p/\log X)}{\pi\sqrt p}\sum_i b_i p^{-ihu_i}
 \quad(Y<p\le X).
\]
Replace $p^{-ix}$ provisionally by independent uniform phases on the unit circle. The summands are centered, bounded uniformly, and independent. By (6)--(7) their total variance is
$\tfrac12\sum_p|d_p|^2=(\sigma_c^2\sum_i b_i^2+o(1))L$.
The bounded summands satisfy Lindeberg's condition. All fixed normalized moments converge to the Gaussian moments as well: the cumulants of each order $k\ge3$ are bounded in total by $O_k(\sum_p p^{-k/2})=O_k(1)$, while the second cumulant is the variance.

For actual height averaging, expand each fixed $k$th moment in prime products. The diagonal products give exactly the independent-phase moments by unique factorization. For unequal products $n,m\le X^k$, the mean of $(n/m)^{ix}$ has magnitude $O(X^k/T)$ since $|\log(n/m)|\gg X^{-k}$. The sum of absolute coefficients of the expanded moment is $O_k((C\sqrt X)^k)$ for fixed $b$ and $c$. The entire off-diagonal error is $T^{-1+o(1)}$. Gaussian moment determinacy and Cram\'er--Wold finish the joint convergence. The same second-moment calculation for one coordinate gives the stated uniform bound. For fixed-degree moments the estimates for these explicit prime sums are uniform under bounded translations of the height interval.
\end{proof}

\section{Mean dispersion about a block median}
For $I_x=(x,x+h]$, set $D_h(U;x)=h^{-1}H_{I_x}(U)$. We claim that
\[
 \frac1T\int_T^{2T-h}D_h(F_0;x)\,dx
       =(\kappa_c+o(1))\sqrt L .                                    \tag{8}
\]
By \eqref{eq:err}, \eqref{eq:slow}, translation invariance and the $L^1$ Lipschitz bound for dispersion, replacing $F_0$ by $Z$ in this average costs $O(1+L^{1/4})$. Indeed one may subtract $B(x)$ throughout $I_x$, then average the error over $0\le v\le h$. The same comparison on the slightly enlarged integration interval is valid.

For $Z$ the choice of center zero gives the upper bound $(\E|G_1|+o(1))\sqrt L$ by Lemma~\ref{le:prime} and uniform integrability. For the matching lower bound, take $s$ independent uniform offsets $U_i\in[0,1]$. The infimum of an expectation is at least the expectation of the infimum, so
\[
 D_h(Z;x)\ge \E_U\inf_m\frac1s\sum_{i=1}^s |Z(x+hU_i)-m|.
\]
For almost every fixed tuple of offsets the entries are distinct. The sample dispersion is continuous in those entries and is bounded above by their mean absolute value. After division by $\sqrt L$, height averaging therefore converges to $\E\inf_m s^{-1}\sum_i|G_i-m|$. The uniform second-moment bound justifies convergence of first moments and integration over the offsets. Trimming a subinterval of length $h\le1$ at the upper endpoint gives an $o(1)$ change in the normalized averages by Cauchy--Schwarz. Finally, as $s\to\infty$, a sample median of the independent Gaussian values converges almost surely to zero. Their sample mean absolute value converges to $\E|G_1|$. The sample dispersion converges in mean to that same number by its uniform second-moment bound. This proves (8), since
\[
 \E|G_1|=\sigma_c\sqrt{2/\pi}=\kappa_c.
\]
The independent values in this argument arise only after removing the low prime frequencies. Centering a block at its first zero-count value would instead introduce a shared fluctuating endpoint.

\section{From local dispersion to the optimal global cost}
To prove the lower bound in Theorem~\ref{th:main}, take $h=\lambda$. Shift a grid of mesh $h$ by a uniform offset in $[0,h]$, and use its complete blocks inside $(T,2T]$. The expected sum of $H_I(F_0)$ is
\[
 h^{-1}\int_T^{2T-h}H_{I_x}(F_0)\,dx.
\]
The left side of (1) is independent of the shift. Equation (8) gives the asserted lower bound.

For the upper bound extend the same shifted grid to the whole real line, and take
$g=\lfloor B(x)+1/2\rfloor$ on a block with left endpoint $x$, restricted to $(T,2T)$. Use $F_0$ and its absolute endpoint values in (2). After averaging over the grid shift,
\[
 \E_{\rm grid}\int_T^{2T}|F_0(t)-g(t)|\,dt
 \le \int_T^{2T} |Z(t)|\,dt+O(T(1+L^{1/4}))
 =(\kappa_c+o(1))T\sqrt L.                                         \tag{9}
\]
Indeed for each $t$ the distance from $t$ to its preceding grid point is uniform on $[0,h]$, so \eqref{eq:slow} bounds the averaged discrepancy of the baselines. Consecutive rounded values differ on average by $O(1)$ by the same estimate, giving expected total interior jump cost $O(\lambda T/h)=O(T)$. More explicitly the expected sum of jump sizes is $h^{-1}$ times the integral of the one-step difference over possible interior grid points. The two endpoint corrections cost $o(T)$: the argument bound following (3), the affine slope, and the elementary absolute bound $|P_{\rm lo}|\ll\sqrt X=T^{o(1)}$ give a uniform endpoint bound $O(\lambda(1+\log T+\sqrt X))$. Some grid has total cost at most the average. Equations (2) and (9) complete Theorem~\ref{th:main}.

For Theorem~\ref{th:full}, set $h=1$. Then $J=F_0(2T)-F_0(T)=\beta_c T+O(\log T)$, and (1), averaged as before, yields
\[
 \cost_\lambda\ge\lambda|J|+(\kappa_c+o(1))T\sqrt{\log\log T}
 \qquad(\lambda\ge1),
\]
because the subtracted $|J|=O(T)$ is lower order. For the upper bound construct a monotone integer step function by rounding $m(t)$ to the nearest integer and clipping the result to the interval between $F_0(T)$ and $F_0(2T)$. The sign of $J$ eventually equals the sign of the fixed nonzero slope $\beta_c$. Thus both endpoint corrections preserve monotonicity and $V(g)=|J|$.

Clipping affects at most $O_c(\log T)$ units of length at the two ends, since the difference between each endpoint count and its affine approximation is $O(\log T)$ and $|\beta_c|>0$. The additional integrated error is $O_c((\log T)^2)$, and rounding costs at most $T/2$. At $h=1$ the low sum has just the primes $p\le e$ and is bounded. Hence
\[
 \int_T^{2T}|F_0-m|\,dt
   =(\kappa_c+o(1))T\sqrt{\log\log T}
\]
by \eqref{eq:err} and the one-point limit in Lemma~\ref{le:prime}. The monotone conclusion of Lemma~\ref{le:flow} supplies an actual matching of the entire smaller list at this cost. Its converse lower bound follows from $\cost_1\le W_T+|J|$. Using that matching in the definition of $\cost_\lambda$ gives the claimed uniform expansion even if $\lambda$ grows arbitrarily quickly.

\section{Equal conductors and the cost of the two endpoints}
Suppose $\beta_c=0$, so $m(x)=m$ is constant. Set $s=\sqrt{\log\log T}$,
$A=F_0(T)$, $B=F_0(2T)$, and let $G$ have the normal law in Lemma~\ref{le:prime}.
For all real $z$ put
\[
 f_T(z)=s\,\E\,|G-(z-m)/s|.
\]
\begin{theorem}[The complete price curve at equal conductors]
Uniformly in $\lambda\ge1$,
\[
 \cost_\lambda=\min_{z\in\mathbb Z}
  \{T f_T(z)+\lambda(|z-A|+|z-B|)\}+o(Ts).
\]
Let $z_*$ be the point of the closed interval with endpoints $A,B$ nearest to $m$. The maximum-cardinality matching costs
\[
 W_T=T f_T(z_*)+o(Ts).
\]
These are absolute errors even along sequences on which $|z_*-m|/s$ tends to infinity.
\end{theorem}
\begin{proof}
This will also prove Theorem~\ref{th:end}. Take $h=1/\log\log T$ and write $\eta_x$ for the distribution of $(F_0(x+hU)-m)/s$ when $U$ is uniform on $[0,1]$. Write $\gamma$ for the law of $G$. We claim
\[
 T^{-1}\int_T^{2T}W_1(\eta_x,\gamma)\,dx=o(1),
\]
where $W_1$ is the Wasserstein distance on the line, equivalently the integral of the absolute difference of the two distribution functions. By the prime approximation, replacing the variable by $Z(x+hU)/s$ costs $o(1)$ in this averaged distance. In fact the $L^1$ approximation error divided by $s$ tends to zero, and the low-prime mean square is $O(1+\log\log Y)=O(1+\log(1/h))=o(s^2)$ by the Dirichlet-polynomial mean-value estimate. Here $L/\log\log T\to1$. For the high prime part and any real $v$, the first two moments over $x$ of the random distribution-function value at $v$ tend to $\gamma(({-}\infty,v])$ and its square. This follows by Fubini and the joint Gaussian limit at two distinct offsets, with the diagonal having zero measure. Therefore the expected absolute distribution-function discrepancy tends to zero at every $v$. Integrate first over $[-R,R]$, then bound the expected integral outside it by $O(1/R)$ using the uniform normalized second-moment bound and the Gaussian second moment. This proves the claim.

Given an optimal partial matching, define the integer step function $g(t)$ by adding the cumulative signed omissions through $t$ to $A$. Its extended endpoints are $A,B$. The number of omissions dominates $V(g)$. Across each cut the remaining difference $F_0-g$ equals the signed number of matched paths crossing it. Integration shows that total movement dominates $\int |F_0-g|$. On the other hand the constant nearest integer to $m$ in (2) gives $\cost_\lambda\ll Ts+\lambda\log T$ by the prime approximation and the endpoint bound. Consequently $V(g)\ll Ts/\lambda+\log T$.

Extend $g$ to the right by $B$. Replacing its transport integral by
\[
 \int_T^{2T}\!\int_0^1 |F_0(x+hu)-g(x)|\,du\,dx
\]
costs $o(Ts)$ uniformly in $\lambda$. Indeed replacing $g(x)$ by $g(x+hu)$ inside costs at most $hV(g)$ after integration. Translating the resulting interval of integration changes its ends at cost $O(h(\log T+V(g)))$, since $|F_0-m|=O(\log T)$ throughout the slightly enlarged dyadic interval and $|g-m|\le |A-m|+V(g)$. The claim on $W_1$ now shows that the last display differs from $\int_T^{2T}f_T(g(x))\,dx$ by $o(Ts)$, uniformly in the values of $g$, since absolute deviation is a 1-Lipschitz test. At each $x$, $V(g)\ge |g(x)-A|+|g(x)-B|$. Integrate this inequality divided by $T$ and combine the costs to obtain the required minimum as a lower bound. Conversely, for any constant integer $g=z$ in (2), its extended variation is exactly $|z-A|+|z-B|$, and its movement integral is $Tf_T(z)+o(Ts)$ uniformly in $z$. For this last assertion use the one-point global $W_1$ convergence, which follows from the one-point high-prime limit, uniform integrability and the same low-prime/error estimates. This proves the first formula.

For $\lambda>T/2$ an optimal matching pairs the entire smaller list, since adding a pair costs at most $T$ and saves $2\lambda$. Projection of an integer $z$ onto the interval between $A,B$ decreases the variation price by $2\lambda$ per step and can increase the movement term by at most $T$ per step. Thus the minimum is attained in the interval. There the variation price is identically $\lambda|B-A|$ and the minimum of the Gaussian absolute deviation is taken at an integer within one of $z_*$, because its median is $m$ in unscaled coordinates. Rounding changes $Tf_T$ by at most $T=o(Ts)$. Subtracting the forced omissions gives the formula for $W_T$.
\end{proof}

\subsection*{A limiting law for a randomly chosen height}
Take $t$ uniformly in $[X_0,2X_0]$ with $X_0\to\infty$ and consider the lists on $(t,2t]$, still with $\beta_c=0$. Let $G',G''$ be independent copies of $G$, let $Z_*$ be the point of their closed interval nearest zero and put $H(v)=\E|G-v|$, a deterministic even function. Then
\[
 \frac{W_t}{t\sqrt{\log\log t}}\ \Longrightarrow\ H(Z_*),\qquad
 \Pr(Z_*=0)=\tfrac12,\qquad
 \Pr(|Z_*|>u)=2\bigl(1-\Phi_0(u/\sigma_c)\bigr)^2\quad(u\ge0),       \tag{10}
\]
where $\Phi_0$ is the standard normal distribution function. The limiting law has an atom of mass exactly one half at its minimum $H(0)=\kappa_c$. The remainder is continuous since $H$ increases strictly on the positive half-line.

To see this, use the prime approximation of Lemma~\ref{le:prime} at $h=1$, with a single cutoff $X=X_0^{1/(\log\log X_0)^{1/4}}$ for both $t$ and $2t$. Both averaged errors divided by $s_{X_0}$ tend to zero by the lemma of Bombieri and Perelli on the two dyadic intervals. The low prime part is bounded. In the phase model the two high-frequency summands for each prime are now proportional to
$\Im(a_p U_p)$ and $\Im(a_p U_p^2)$, where $U_p$ is the same uniform unit phase. They have covariance zero by orthogonality of the first and second harmonics, and are independent from the summands at all other primes. The bounded-summand cumulant proof gives a joint limit of two independent Gaussians with variance $\sigma_c^2$. The height-moment comparison remains valid: prime products in a moment of fixed order $k$ are now at most $X^{2k}=X_0^{o(1)}$, and their equality imposes exactly the same primewise exponent conditions as phase averaging. Thus the normalized endpoint pair converges to $(G',G'')$. Replacing $s_{X_0}$ by $s_t$ changes the normalization by a factor tending uniformly to one. Theorem~\ref{th:end} and continuity give (10). The formula for the tail follows because both endpoints must exceed $u$ or both must be below $-u$.

\section{Interpretation and open questions}
The norm in the leading constant comes from the mean square on primes of $\sum c_j\chi_j(p)$. It detects common primitive factors with multiplicity. For instance, changing a packet $(1,1,0)$ to $(1,0,1)$ has the same leading cost as comparing its two distinct factors directly. Changing $(2,0,0)$ to $(0,1,1)$ multiplies that cost by $\sqrt3$. This conclusion concerns a minimum over all admissible matchings, including matchings that exchange the labels of shared factors.

\begin{corollary}[A hard distance cap]
Let $D_T(K)$ be the fewest omitted occurrences when each pair has distance at most $K/\log T$. For $K\to\infty$ with $K\le\log T$,
\[
 \frac{D_T(K)}{T\log T}=o\!\left(\frac{\sqrt{\log K}}{K}\right).
\]
\end{corollary}
\begin{proof}
Put $\lambda=K/(2\log T)$. In an optimal matching at price $\lambda$ no pair can have length greater than $2\lambda$, since deleting it would be cheaper. If its omission count is $D$, using this matching at price $\lambda/2$ gives
$(\lambda/2)D\le\cost_\lambda-\cost_{\lambda/2}$. Theorem~\ref{th:main} gives both costs the same leading term $\kappa_c T\sqrt{\log K}$. Their difference is therefore of smaller order, proving the claim.
\end{proof}

\paragraph{Scope of the distinct-zeros problem.}
Bombieri and Perelli \cite[p.~274]{bp} discuss the expectation that almost all zeros of two distinct primitive functions are distinct. A distance asymptotic after optimizing the tradeoff between motion and omissions does not determine the density of exact common zeros. The companion note studies the problem in which a strict distance cap replaces the deletion price. The dyadic averaging here also leaves open local uniform statements in the short windows treated by ProveIt.

\paragraph{Further research questions.}
At fixed $M$ one can ask for a limit of $\cost_{M/\log T}/T$ and for its dependence on $M$. Finer joint statistics than the mesoscopic Gaussian input may be needed. For equal conductor products the random-height law (10) gives an atom of mass one half. Obtaining a quantitative convergence rate, especially near that atom, would require sharper control of the empirical distributions on unit intervals. Another question is whether the $\kappa_c$ law holds uniformly over conductors growing with $T$.

\begin{thebibliography}{9}
\bibitem{bp} E. Bombieri and A. Perelli, \emph{Distinct zeros of $L$-functions},
Acta Arith. 83 (1998), 271--281.
\href{https://doi.org/10.4064/aa-83-3-271-281}{doi:10.4064/aa-83-3-271-281}.
\bibitem{bourgade} P. Bourgade, \emph{Mesoscopic fluctuations of the zeta zeros},
Probab. Theory Relat. Fields 148 (2010), 479--500.
\href{https://arxiv.org/abs/0902.1757}{arXiv:0902.1757}.
\bibitem{fujii} A. Fujii, \emph{On the zeros of Dirichlet $L$-functions (V)},
Acta Arith. 28 (1976), 395--403.
\href{https://matwbn.icm.edu.pl/ksiazki/aa/aa28/aa2847.pdf}{Publisher's text}.
\bibitem{repo} V. Reshetnikov, \emph{Beyond the Riemann hypothesis for Dirichlet
$L$-functions: consequences, transfers and limits}, ProveIt, snapshot
\texttt{cc919c585f90}, 27 November 2026.
\href{https://github.com/VladimirReshetnikov/ProveIt/blob/cc919c585f90be8e0fba7a1b90c70afdea39998e/Analysis/GeneralizedRiemannHypothesis/Article/sections/rec.tex}{Section on recovering zeros and characters},
especially the corollaries on local matching and moving or deleting zeros.
\end{thebibliography}
\end{document}
